%%% Внесите свои данные - Input your data
%%
%%
\newcommand{\Author}{М.Р.\,Сабирзянов} % И.О. Фамилия автора
\newcommand{\AuthorFull}{Сабирзянов Мирсаит Раисович} % Фамилия Имя Отчество автора
\newcommand{\AuthorFullDat}{Сабирзянову Мирсаиту Раисовичу} % Фамилия Имя Отчество автора в дательном падеже (Кому? Студенту...)
\newcommand{\AuthorFullVin}{Сабирзянова Мирсаита Раисовича} % в винительном падеже (Кого? что?  Програмиста ...)
\newcommand{\AuthorPhone}{} % номер телефорна автора для оперативной связи
\newcommand{\Supervisor}{А.В.\,Сергеев} % И. О. Фамилия научного руководителя
\newcommand{\SupervisorFull}{Сергеев Анатолий Васильевич} % Фамилия Имя Отчество научного руководителя
\newcommand{\SupervisorVin}{А.В.\,Сергеева} % И. О. Фамилия научного руководителя  в винительном падеже (Кого? что? Руководителя ...)
\newcommand{\SupervisorJob}{к.т.н., доцент ВШПИ} %
\newcommand{\SupervisorJobVin}{доцента} % в винительном падеже (Кого? что?  Програмиста ...)
\newcommand{\SupervisorDegree}{} %
\newcommand{\SupervisorTitle}{} % 
%%
%%
%Руководитель, УТВЕРЖДАЮЩИЙ задание
\newcommand{\Head}{А.В.\,Щукин} % И. О. Фамилия руководителя подразделения (руководителя ОП)
\newcommand{\HeadDegree}{Руководитель ОП}% Только должность:   
%Руководитель ОП % <- ПРАВИЛЬНЫЙ ВАРИАНТ ДЛЯ 09.03.03 и 09.04.03
%Заведующий % кафедрой
%Директор % Высшей школы
%Зам. директора
\newcommand{\HeadDep}{} % заменить на краткую аббревиатуру подразделения или ОСТАВИТЬ ПУСТЫМ, если утверждает руководитель ОП

%%% Руководитель, принимающий заявление
\newcommand{\HeadAp}{И.О.\,Фамилия} % И. О. Фамилия руководителя подразделения (руководителя ОП)
\newcommand{\HeadApDegree}{Руководитель ОП}% Только должность:   
%Руководитель ОП 
%Заведующий кафедрой
%Директор Высшей школы
\newcommand{\HeadApDep}{} % заменить на краткую аббревиатуру подразделения или оставить пустым, если утверждает руководитель ОП
%%% Консультант по нормоконтролю
\newcommand{\ConsultantNorm}{Е.Е.\,Андрианова} % И. О. Фамилия консультанта по нормоконтролю. ТОЛЬКО из числа ППС!
\newcommand{\ConsultantNormDegree}{Старший преподаватель ВШПИ} %
%%% Первый консультант
\newcommand{\ConsultantExtraFull}{Фамилия Имя Отчетство} % Фамилия Имя Отчетство дополнительного консультанта 
\newcommand{\ConsultantExtra}{И.О.\,Фамилия} % И. О. Фамилия дополнительного консультанта 
\newcommand{\ConsultantExtraDegree}{должность, степень} % 
\newcommand{\ConsultantExtraVin}{И.О.\,Фамилию} % И. О. Фамилия дополнительного консультанта в винительном падеже (Кого? что? Руководителя ...)
\newcommand{\ConsultantExtraDegreeVin}{должность, степень} %  в винительном падеже (Кого? что? Руководителя ...)
%%% Второй консультант
\newcommand{\ConsultantExtraTwoFull}{Фамилия Имя Отчетство} % Фамилия Имя Отчетство дополнительного консультанта 
\newcommand{\ConsultantExtraTwo}{И.О.\,Фамилия} % И. О. Фамилия дополнительного консультанта 
\newcommand{\ConsultantExtraTwoDegree}{должность, степень} % 
\newcommand{\ConsultantExtraTwoVin}{И.О.\,Фамилию} % И. О. Фамилия дополнительного консультанта в винительном падеже (Кого? что? Руководителя ...)
\newcommand{\ConsultantExtraTwoDegreeVin}{должность, степень} %  в винительном падеже (Кого? что? Руководителя ...)
\newcommand{\Reviewer}{И.О.\,Фамилия} % И. О. Фамилия резензента. Обязателен только для магистров.
\newcommand{\ReviewerDegree}{должность, степень} % 
%%
%%
\renewcommand{\thesisTitle}{Проектирование и реализация iOS-приложения как комплексного помощника для мусульман}
\newcommand{\thesisDegree}{работа бакалавра}% магистерская диссертация %c 2020, И ТОЛЬКО для специалитата дипломный проект и дипломная работа
\newcommand{\thesisTitleEn}{Design and implementation of an iOS application as a comprehensive assistant for Muslims} %2020
\newcommand{\thesisDeadline}{18.05.2026} %Последний день преддипломной практики согласно учебному плану.
\newcommand{\thesisStartDate}{04.05.2026}
\newcommand{\thesisYear}{2026} % заменить на год защиты
\newcommand{\approveYear}{202\underline{\hspace*{0.01\textheight}}} % <- НЕ МЕНЯТЬ, ТОЛЬКО С 2030го :)
%%
%%
\newcommand{\group}{5130903/20301} % заменить вместо N номер группы
\newcommand{\thesisSpecialtyCode}{09.03.03}% код направления подготовки
\newcommand{\thesisSpecialtyTitle}{Прикладная информатика} % наименование направления/специальности
\newcommand{\thesisOPPostfix}{03} % последние цифры кода образовательной программы (после <<_>>)
\newcommand{\thesisOPTitle}{Интеллектуальные инфокоммуникационные технологии}% наименование образовательной программы
%%
%%
\newcommand{\institute}{
Институт компьютерных наук и~кибербезопасности
%Институт компьютерных наук и~кибербезопасности
%Гуманитарный институт
%Инженерно-строительный институт
%Институт биомедицинских систем и технологий
%Институт металлургии, машиностроения и транспорта
%Институт передовых производственных технологий
%Институт прикладной математики и механики
%Институт физики, нанотехнологий и телекоммуникаций
%Институт физической культуры, спорта и туризма
%Институт энергетики и транспортных систем
%Институт промышленного менеджмента, экономики и торговли
}%
%%
%%




%%% Задание ключевых слов и аннотации
%%
%%
%% Ключевых слов от 3 до 5 слов или словосочетаний в именительном падеже именительном падеже множественного числа (или в единственном числе, если нет другой формы) по правилам русского языка!!!
%%
%%
\newcommand{\keywordsRu}{iOS-приложение, SwiftUI, RAG, искусственный интеллект, FastAPI, Spring Boot, OCR, халяль, Коран, LLM} % ВВЕДИТЕ ключевые слова по-русски
%%
%%
\newcommand{\keywordsEn}{iOS application, SwiftUI, RAG, artificial intelligence, FastAPI, Spring Boot, OCR, halal, Quran, LLM} % ВВЕДИТЕ ключевые слова по-английски
%%
%%
%% Реферат ОТ 1000 ДО 1500 знаков на русский или английский текст
%%
%Реферат должен содержать:
%- предмет, тему, цель ВКР;
%- метод или методологию проведения ВКР:
%- результаты ВКР:
%- область применения результатов ВКР;
%- выводы.

\newcommand{\abstractRu}{%
В выпускной квалификационной работе рассмотрены проектирование, разработка и тестирование многофункционального iOS-приложения для мусульман, объединяющего инструменты анализа халяльности продуктов, чат с искусственным интеллектом, модуль чтения Корана и систему расчета времени намаза.
\texorpdfstring{\par}{}
В работе проведен анализ существующих исламских цифровых сервисов и сформулированы функциональные и нефункциональные требования к системе. Разработана архитектура программного комплекса, включающая iOS-клиент на SwiftUI, backend-сервис на Spring Boot и отдельный LLM-сервис на FastAPI. Для формирования ответов AI-ассистента реализован retrieval-augmented generation (RAG) подход с использованием векторного поиска по текстам Корана.
\texorpdfstring{\par}{}
Реализован модуль анализа состава продуктов с использованием OCR и классификации ингредиентов по категориям «халяль», «харам» и «сомнительный». Выполнено проектирование REST API, базы данных PostgreSQL и механизма взаимодействия между сервисами.
\texorpdfstring{\par}{}
Проведено тестирование клиентской и серверной частей системы, а также оценка качества работы RAG-подсистемы и проверки состава продуктов. Результаты показали корректность работы основных модулей приложения и возможность применения разработанной системы в качестве комплексного цифрового помощника для мусульман.%
} % ВВЕДИТЕ текст аннотации по-русски
%%
%%
\newcommand{\abstractEn}{%
This thesis presents the design, development and testing of a multifunctional iOS application for Muslims that combines halal product analysis, an AI-based assistant, Quran reading functionality and prayer time calculations.
\texorpdfstring{\par}{}
The paper includes an analysis of existing Islamic digital services and defines functional and non-functional requirements for the system. A software architecture consisting of an iOS client developed with SwiftUI, a Spring Boot backend and a separate FastAPI-based LLM service was designed. A retrieval-augmented generation (RAG) approach with vector search over Quran texts was implemented for the AI assistant.
\texorpdfstring{\par}{}
A product ingredient analysis module based on OCR and ingredient classification into halal, haram and doubtful categories was implemented. REST API architecture, PostgreSQL database structure and service interaction mechanisms were also designed.
\texorpdfstring{\par}{}
Testing of the client, backend and LLM subsystems was carried out, including evaluation of the RAG pipeline and ingredient recognition quality. The results demonstrated the correctness of the implemented modules and the applicability of the developed system as a comprehensive digital assistant for Muslims.%
} % ВВЕДИТЕ текст аннотации по-английски


%%% РАЗДЕЛ ДЛЯ ОФОРМЛЕНИЯ ПРАКТИКИ
%Место прохождения практики
\newcommand{\PracticeType}{Отчет о прохождении производственной (научно-исследовательской работы) практики}

\newcommand{\Workplace}{ФГАОУ ВО «СПбПУ», ИКНК, ВШ ПИ} % TODO Rename this variable

% Даты начала/окончания
\newcommand{\PracticeStartDate}{04.05.2026}%
\newcommand{\PracticeEndDate}{18.05.2026}%
%%

\newcommand{\School}{Высшая школа программной инженерии}
\newcommand{\practiceTitle}{Проектирование и реализация iOS-приложения как комплексного помощника для мусульман}


%% ВНИМАНИЕ! Необходимо либо заменить текст аннотации (ключевых слов) на русском и английском, либо удалить там весь текст, иначе в свойства pdf-отчета по практике пойдет шаблонный текст.

%%% Не меняем дальнейшую часть - Do not modify the rest part
%%
%%
%%
%%
\ifnumequal{\value{docType}}{1}{% Если ВКР, то...
	\newcommand{\DocType}{Выпускная квалификационная работа}
	\newcommand{\pdfDocType}{\DocType~(\thesisDegree)} %задаём метаданные pdf файла
	\newcommand{\pdfTitle}{\thesisTitle}
}{% Иначе 
	\newcommand{\DocType}{\PracticeType}
	\newcommand{\pdfDocType}{\DocType} %задаём метаданные pdf файла
	\newcommand{\pdfTitle}{\practiceTitle}
}%
\newcommand{\HeadTitle}{\HeadDegree~\HeadDep}
\newcommand{\HeadApTitle}{\HeadApDegree~\HeadApDep}
\newcommand{\thesisOPCode}{\thesisSpecialtyCode\_\thesisOPPostfix}% код образовательной программы
\newcommand{\thesisSpecialtyCodeAndTitle}{\thesisSpecialtyCode~\thesisSpecialtyTitle}% Код и наименование направления/специальности
\newcommand{\thesisOPCodeAndTitle}{\thesisOPCode~\thesisOPTitle} % код и наименование образовательной программы
%%
%%
\hypersetup{%часть болка hypesetup в style
		pdftitle={\pdfTitle},    % Заголовок pdf-файла
		pdfauthor={\AuthorFull},    % Автор
		pdfsubject={\pdfDocType. Шифр и наименование направления подготовки: \thesisSpecialtyCodeAndTitle. \abstractRu},      % Тема
		pdfcreator={LaTeX, SPbPU-student-thesis-template},     % Приложение-создатель
%		pdfproducer={},  % Производитель, Производитель PDF % будет выставлена автоматически
		pdfkeywords={\keywordsRu}
}
%%
%%
%% вспомогательные команды
\newcommand{\firef}[1]{рис.\ref{#1}} %figure reference
\newcommand{\taref}[1]{табл.\ref{#1}}	%table reference
%%
%%
%% Архивный вариант задания ключевых слов, аннотации и благодарностей 
% Too hard to export data from the environment to pdf-info
% https://tex.stackexchange.com/questions/184503/collecting-contents-of-environment-and-store-them-for-later-retrieval
%заменить NewEnviron на newenvironment для распознавания команды в TexStudio
%\NewEnviron{keywordsRu}{\noindent\MakeUppercase{\BODY}}
%\NewEnviron{keywordsEn}{\noindent\MakeUppercase{\BODY}}
%\newenvironment{abstractRu}{}{}
%\newenvironment{abstractEn}{}{}
%\newenvironment{acknowledgementsRu}{\par{\normalfont \acknowledgements.}}{}
%\newenvironment{acknowledgementsEn}{\par{\normalfont \acknowledgementsENG.}}{}


%%% Переопределение именований %%% Не меняем - Do not modify
%\newcommand{\Ministry}{Минобрнауки России} 
\newcommand{\Ministry}{Министерство науки и высшего образования Российской~Федерации} %с 2020
\newcommand{\SPbPU}{Санкт-Петербургский политехнический университет Петра~Великого}
\newcommand{\SPbPUOfficialPrefix}{Федеральное государственное автономное образовательное учреждение высшего образования}
\newcommand{\SPbPUOfficialShort}{ФГАОУ~ВО~<<СПбПУ>>}
%% Пробел между И. О. не допускается.
\renewcommand{\alsoname}{см. также}
\renewcommand{\seename}{см.}
\renewcommand{\headtoname}{вх.}
\renewcommand{\ccname}{исх.}
\renewcommand{\enclname}{вкл.}
\renewcommand{\pagename}{Pages}
\renewcommand{\partname}{Часть}
\renewcommand{\abstractname}{\textbf{Аннотация}}
\newcommand{\abstractnameENG}{\textbf{Annotation}}
\newcommand{\keywords}{\textbf{Ключевые слова}}
\newcommand{\keywordsENG}{\textbf{Keywords}}
\newcommand{\acknowledgements}{\textbf{Благодарности}}
\newcommand{\acknowledgementsENG}{\textbf{Acknowledgements}}
\renewcommand{\contentsname}{Content} % 
%\renewcommand{\contentsname}{Содержание} % (ГОСТ Р 7.0.11-2011, 4)
%\renewcommand{\contentsname}{Оглавление} % (ГОСТ Р 7.0.11-2011, 4)
\renewcommand{\figurename}{Рис.} % Стиль СПбПУ
%\renewcommand{\figurename}{Рисунок} % (ГОСТ Р 7.0.11-2011, 5.3.9)
\renewcommand{\tablename}{Таблица} % (ГОСТ Р 7.0.11-2011, 5.3.10)
%\renewcommand{\indexname}{Предметный указатель}
\renewcommand{\listfigurename}{Список рисунков}
\renewcommand{\listtablename}{Список таблиц}
\renewcommand{\refname}{\fullbibtitle}
\renewcommand{\bibname}{\fullbibtitle}

\newcommand{\chapterEnTitle}{Сhapter title} % <- input the English title here (only once!) 
\newcommand{\chapterRuTitle}{Название главы}          % <- введите 
\newcommand{\sectionEnTitle}{Section title} %<- input subparagraph title in english
\newcommand{\sectionRuTitle}{Название подраздела} % <- введите название подраздела по-русски
\newcommand{\subsectionEnTitle}{Subsection title} % - input subsection title in english
\newcommand{\subsectionRuTitle}{Название параграфа} % <- введите название параграфа по-русски
\newcommand{\subsubsectionEnTitle}{Subsubsection title} % <- input subparagraph title in english
\newcommand{\subsubsectionRuTitle}{Название подпараграфа} % <- введите название подпараграфа по-русски